%% toc: P3 27 27.6,27.7,27.8
\section{Keisler's and Shelah's isomorphic-ultrapowers theorems and their hypotheses}
\label{sec:p03-w03-p03-018-shelah-ultrapowers}

\Cref{thm:p03-keisler-shelah} says that two structures are elementarily
equivalent exactly when some ultrapower of one is isomorphic to some
ultrapower of the other.  That existential form hides what the original
theorem fixes in advance: one infinite index cardinal \(\lambda\), a cardinal
bound \(\mu\) computed from \(\lambda\), and a single ultrafilter \(D\) over
\(\lambda\) that works simultaneously for every pair of elementarily
equivalent models below the bound.  This section states the theorem in the
form Shelah published it \cite[224]{Shelah1971Ultrapowers}; states and proves
Keisler's earlier theorems, which need an instance of the continuum
hypothesis, from his article \cite{Keisler1961ElementaryClasses}, and
compares the two sets of hypotheses; derives the explicit index sets behind
the existential form; and separates three relations that are easily
conflated: elementary equivalence of two structures, isomorphism of their
ultrapowers, and isomorphism of the structures themselves.

The inputs are an infinite cardinal \(\lambda\) and first-order models of a
common language.  The conclusions are an ultrafilter over \(\lambda\), an
isomorphism between ultrapowers of elementarily equivalent models of small
power, and, as corollaries, explicit index sets for the existential form.
Ultraproducts, ultrapowers and \L o\'s's theorem are as in
\cref{sec:types-ultraproducts-05}; the ambient set theory includes the Axiom
of Choice, and \cref{def:p03-ultrafilter-choice-boundary} records which
steps in this part use it.

\begin{hbconvention}{def:p03-w03-p03-018-notation}{Shelah's notation}
Except in \cref{sec:p03-w03-p03-018-keisler}, the notation in this section
is Shelah's
\cite[225]{Shelah1971Ultrapowers}.  The letter \(\lambda\) is a fixed
infinite cardinal and
\[
  \mu=\min\{\mu:\lambda^\mu>\lambda\},
\]
the least cardinal \(\mu\) for which the power \(\lambda^\mu\) exceeds
\(\lambda\); the bound variable and the defined cardinal share the letter, as
in the source.  The exponent base is \(\lambda\) itself, not \(2\): the bound
is defined by \(\lambda^\mu>\lambda\) both in the statement of the theorem
\cite[224]{Shelah1971Ultrapowers} and in the notation paragraph.
Models are written \(M\), \(N\).  The source writes the universe of \(M\)
between single bars; here it is \(\dom M\), because single bars denote
cardinality in this handbook.  The cardinality of the universe, the
\emph{power} of \(M\), is \(\lVert M\rVert\), as in the source.  An
ultrafilter \(D\) \emph{over} \(\lambda\) is an ultrafilter on the set
\(\lambda=\{i:i<\lambda\}\), and \(M^\lambda/D\) is the ultrapower of \(M\)
with index set \(\lambda\) in the sense of \cref{sec:types-ultraproducts-05},
where it is written \(\mathcal M^I/\mathcal U\).  Elementary equivalence is
\(M\equiv N\) and isomorphism is \(M\cong N\).  All cardinal arithmetic is
computed in the ambient set theory with Choice, using
\cref{thm:p06-cardinal-arithmetic-laws}.
\end{hbconvention}

The cardinal \(\mu\) controls which models the theorem covers, so its basic
behaviour comes first. Its regularity is proved separately in
\cref{thm:p03-w03-p03-017-mu-regular}.

\begin{hblemma}{thm:p03-w03-p03-018-mu}{proved=annex:p03-w03-p03-018-mu}{The bound \(\mu\) attached to \(\lambda\)}
Let \(\lambda\) be an infinite cardinal and \(\mu\) as in
\cref{def:p03-w03-p03-018-notation}.
\begin{enumerate}
  \item The cardinal \(\mu\) exists, is infinite, and satisfies
        \(\mu\le\lambda\).
  \item \(\lambda^\nu=\lambda\) for every cardinal \(\nu\) with
        \(1\le\nu<\mu\).
  \item If \(\lambda=\aleph_0\), then \(\mu=\aleph_0\).
  \item If \(\lambda=2^\kappa\) for an infinite cardinal \(\kappa\), then
        \(\kappa<\mu\).  In particular \(\aleph_0<\mu\) when
        \(\lambda=2^{\aleph_0}\), and \(2^{\aleph_0}<\mu\) when
        \(\lambda=2^{2^{\aleph_0}}\).
\end{enumerate}
\end{hblemma}

\subsection{The theorem as published}

\begin{hbtheorem}{thm:p03-w03-p03-018-shelah}{cited=Shelah1971Ultrapowers}{Shelah's isomorphic-ultrapowers theorem}
Let \(\lambda\) be an infinite cardinal and
\(\mu=\min\{\mu:\lambda^\mu>\lambda\}\).  Then there is an ultrafilter \(D\)
over \(\lambda\) such that:
\begin{enumerate}
  \item if \(M\) and \(N\) are elementarily equivalent models of power
        \({<}\,\mu\), then \(M^\lambda/D\) and \(N^\lambda/D\) are
        isomorphic;
  \item if \(M\) is a model of power \({<}\,\mu\) and \(2^\kappa\le2^\lambda\),
        then \(M^\lambda/D\) is \(\kappa^+\)-saturated;
  \item if \(M_k\) and \(N_k\), for \(k<\lambda\), are models of the same
        language, each of cardinality at most \(\chi\) for a fixed cardinal
        \(\chi<\mu\), and the ultraproducts \(\prod_{k<\lambda}M_k/D\) and
        \(\prod_{k<\lambda}N_k/D\) are elementarily equivalent, then they are
        isomorphic.
\end{enumerate}
\end{hbtheorem}

Parts~1 and~2 are proved in
\cref{thm:p03-w03-p03-017-shelah-proof}; part~3 retains the cited
standing recorded below. The statement is the theorem on the first page
of the paper \cite[224]{Shelah1971Ultrapowers}, where \(\lambda\) is introduced as any
cardinality; the standing hypothesis that \(\lambda\) is infinite is fixed in
the source's notation paragraph \cite[225]{Shelah1971Ultrapowers}.  Part~1 is
the precise content of the forward implication of
\cref{thm:p03-keisler-shelah}: the existential phrase ``some ultrapower of
\(M\) is isomorphic to some ultrapower of \(N\)'' is realized with the index
set \(\lambda\) for both factors and with one ultrafilter \(D\) for both,
provided both models have power below \(\mu\).  Four features of the
hypotheses carry the content.

\begin{itemize}
  \item \emph{Index set.}  The index set is the cardinal \(\lambda\) itself,
        and both ultrapowers in part~1 use that same index set.
  \item \emph{Ultrafilter.}  One ultrafilter \(D\) is chosen before any model
        is mentioned, and it serves every pair in part~1 and every model in
        part~2 at once.  The proof lists all pairs of elementarily
        equivalent models whose universes are cardinals below \(\mu\) and
        whose language lies in a fixed language of cardinality \(\lambda\),
        and builds a single \(D\) for the whole list
        \cite[230]{Shelah1971Ultrapowers}.  The theorem gives no explicit
        description of \(D\); by
        \cref{thm:p03-w03-p03-018-comparison} it cannot be principal as soon
        as two non-isomorphic models fall under part~1.
  \item \emph{Cardinal bound.}  Only models of power less than \(\mu\) are
        covered, and \(\mu\le\lambda\) by
        \cref{thm:p03-w03-p03-018-mu}.  For \(\lambda=\aleph_0\) the bound
        admits only finite models; covering countable models needs
        \(\lambda^{\aleph_0}=\lambda\), for instance
        \(\lambda=2^{\aleph_0}\).
  \item \emph{Language.}  The models in part~1 have a common first-order
        language, because they are elementarily equivalent; the statement
        imposes no bound on its cardinality.  The proof assumes without loss
        of generality that a model \(M\) of power below \(\mu\) has a
        language of cardinality at most
        \(2^{\lVert M\rVert}\le\lambda^{\lVert M\rVert}=\lambda\)
        \cite[230]{Shelah1971Ultrapowers}.
\end{itemize}

Saturation in part~2 is not redefined in the source; the handbook's
definition of \(\kappa\)-saturation is the one preceding
\cref{thm:saturated-extension}.  The source states part~3 but omits its
proof (remark~2 at the start of the proof of the main theorem
\cite[230]{Shelah1971Ultrapowers}), so part~3 has a weaker standing in the
source than parts~1 and~2.  The same page notes that, by part~2 together with
an earlier result of Keisler, the ultrafilter \(D\) is \(\lambda^+\)-good.

\subsection{Keisler's conditional theorems of 1961}
\label{sec:p03-w03-p03-018-keisler}

The first proof that elementarily equivalent structures have isomorphic
ultrapowers is Keisler's \cite{Keisler1961ElementaryClasses}.  It is
conditional on cardinal arithmetic, it uses two ultrafilters that are built
after the structures are given, and its index set has the cardinality of the
cardinal \(\alpha\) at which the arithmetic is assumed.  This subsection states
Keisler's hypotheses in his own terms, proves his theorems from them, and
then compares them with \cref{thm:p03-w03-p03-018-shelah}.

\begin{hbconvention}{def:p03-w03-p03-018-keisler-notation}{Keisler's notation}
In this subsection the notation is Keisler's
\cite[479--483]{Keisler1961ElementaryClasses}, with two changes recorded
below.  The Axiom of Choice is assumed throughout, together with the
well-ordering principle and the maximal principle
\cite[479]{Keisler1961ElementaryClasses}.  Cardinals are initial ordinals;
\(\alpha\) and \(\beta\) denote cardinals, \(\alpha^+\) is the least cardinal
above \(\alpha\), and \(\omega\) is the least infinite ordinal.  For a set
\(X\), \(\varkappa(X)\) is its cardinality and \(Y^X\) is the set of functions
from \(X\) to \(Y\); thus \(\varkappa(2^\alpha)\) is the cardinal written
\(2^\alpha\) elsewhere in this handbook.  For a cardinal \(\alpha\),
\(S_\alpha(X)\) is the set of subsets of \(X\) of cardinality less than
\(\alpha\), and
\[
  \overline S_\alpha(X)=\{X-Y:Y\in S_\alpha(X)\},
\]
the set of subsets of \(X\) whose complement in \(X\) has fewer than \(\alpha\)
elements \cite[480]{Keisler1961ElementaryClasses}.  A filter on \(X\) is
proper, and an ultrafilter on \(X\) is a filter containing \(Y\) or \(X-Y\)
for each \(Y\subseteq X\); principal ultrafilters are not excluded.

Structures are relational systems
\(\mathfrak A=\langle A,R_\xi\rangle_{\xi<\varrho}\) and
\(\mathfrak B=\langle B,S_\xi\rangle_{\xi<\varrho}\) of one similarity type:
\(A\) is a nonempty set and each \(R_\xi\) is a relation on \(A\) whose finite
number of argument places is fixed by the type.  The ordinal \(\varrho\)
indexes the predicate symbols, so \(\varkappa(\varrho)\) is their number.  The
language is first-order with identity and with predicate symbols only;
operations are represented by their graphs
\cite[481]{Keisler1961ElementaryClasses}.  \(\mathfrak A^*\) is the set of
sentences true in \(\mathfrak A\), so \(\mathfrak A\equiv\mathfrak B\) exactly
when \(\mathfrak A^*=\mathfrak B^*\).  The ultraproduct
\(\prod_{i\in I}\mathfrak A_i/D\) and the ultrapower \(\mathfrak A^I/D\) are
those of \cref{sec:types-ultraproducts-05}.  \(\mathrm{Pow}(\mathfrak A)\) is
the class of structures isomorphic to \(\mathfrak A^I/D\) for some set \(I\)
and some ultrafilter \(D\) on \(I\) \cite[483]{Keisler1961ElementaryClasses}.
The generalized continuum hypothesis is the statement that
\(\alpha^+=\varkappa(2^\alpha)\) for every infinite cardinal \(\alpha\)
\cite[480]{Keisler1961ElementaryClasses}.

The two changes are these.  Keisler names the similarity type by the letter
\(\mu\) and indexes its predicates by \(\lambda<\varrho\); because \(\lambda\)
and \(\mu\) are Shelah's cardinals in this section, the type is left unnamed
and its predicates are indexed by \(\xi\).  Keisler writes \(0\) for the empty
set; here it is \(\varnothing\).
\end{hbconvention}

\begin{hblemma}{thm:p03-w03-p03-018-keisler-disjoint}{proved=annex:p03-w03-p03-018-keisler-disjoint}{Keisler's disjoint refinement}
Let \(\alpha\) be an infinite cardinal and let
\(\langle X_\xi\rangle_{\xi<\alpha}\) be sets with
\(\varkappa(X_\xi)=\alpha\) for each \(\xi<\alpha\).  Then there is a
sequence \(\langle Y_\xi\rangle_{\xi<\alpha}\) such that
\(Y_\eta\cap Y_\xi=\varnothing\) whenever \(\eta<\xi<\alpha\), and
\(Y_\eta\subseteq X_\eta\) and \(\varkappa(Y_\eta)=\alpha\) for every
\(\eta<\alpha\).
\end{hblemma}

\begin{hbtheorem}{thm:p03-w03-p03-018-keisler-products}{proved=annex:p03-w03-p03-018-keisler-products}{Keisler's ultraproduct theorem}
Let \(\alpha\) be a cardinal with \(\alpha\ge\omega\),
\(\alpha\ge\varkappa(\varrho)\) and \(\varkappa(2^\alpha)=\alpha^+\).  Let
\(\varkappa(I)=\alpha\), and let \(\langle\mathfrak A_i\rangle_{i\in I}\) and
\(\langle\mathfrak B_i\rangle_{i\in I}\) be systems of structures of the type,
each of cardinality at most \(\varkappa(2^\alpha)\).  The following are
equivalent.
\begin{enumerate}
  \item[(i)] There are ultrafilters \(D\), \(E\) on \(I\) with
        \(\overline S_\alpha(I)\subseteq D\), \(\overline S_\alpha(I)\subseteq E\)
        and \(\prod_{i\in I}\mathfrak A_i/D\cong\prod_{i\in I}\mathfrak B_i/E\).
  \item[(ii)] There is no sentence \(\Phi\) with
        \(\{i\in I:\Phi\in\mathfrak A_i^*\}\in\overline S_\alpha(I)\) and
        \(\{i\in I:(\neg\Phi)\in\mathfrak B_i^*\}\in\overline S_\alpha(I)\).
\end{enumerate}
\end{hbtheorem}

This is Theorem~2.2 of the source \cite[484]{Keisler1961ElementaryClasses}, with Lemma~2.1 before it; the proof
\cite[484--486]{Keisler1961ElementaryClasses} is a transfinite back-and-forth
of length \(\alpha^+\) that adapts Cantor's argument for countable dense
orders \cite[486, footnote~3]{Keisler1961ElementaryClasses}.  It enumerates
the products \(\prod_iA_i\) and \(\prod_iB_i\), which have at most
\(\alpha^+\) elements because each factor has at most
\(\varkappa(2^\alpha)=\alpha^+\) elements and \(\varkappa(I)=\alpha\).  It then
extends a pair of sequences of elements one step at a time, preserving
condition~(ii) for the language with a new unary predicate naming each chosen
element, and only at the end builds \(E\) and then \(D\) from the completed
sequences.  So the ultrafilters depend on the two systems.  When every
\(A_i\) has more than one element, \(D\) is determined by \(E\) and the two
sequences \cite[486]{Keisler1961ElementaryClasses}.

\begin{hbtheorem}{thm:p03-w03-p03-018-keisler-powers}{proved=annex:p03-w03-p03-018-keisler-powers}{Keisler's ultrapower theorem}
Let \(\alpha\) be a cardinal with \(\alpha\ge\omega\),
\(\alpha\ge\varkappa(\varrho)\) and \(\varkappa(2^\alpha)=\alpha^+\).  Let
\(\mathfrak A\), \(\mathfrak B\) be structures of cardinality at most
\(\varkappa(2^\alpha)\).  The following are equivalent.
\begin{enumerate}
  \item[(i)] There exist ultrafilters \(D\), \(E\) on \(\alpha\) such that
        \(\mathfrak A^\alpha/D\cong\mathfrak B^\alpha/E\).
  \item[(ii)] \(\mathfrak A^*=\mathfrak B^*\).
\end{enumerate}
Moreover, when (ii) holds, \(D\) and \(E\) in (i) can be chosen to contain
\(\overline S_\alpha(\alpha)\).
\end{hbtheorem}

The equivalence is Theorem~2.3 of the source \cite[486]{Keisler1961ElementaryClasses}, where it is
derived as the case \(I=\alpha\), \(\mathfrak A_i=\mathfrak A\),
\(\mathfrak B_i=\mathfrak B\) of Keisler's Theorem~2.2
\cite[487]{Keisler1961ElementaryClasses}; the final sentence records what
that derivation gives and is proved with it.  Only the single instance
\(\varkappa(2^\alpha)=\alpha^+\) of the continuum hypothesis is used, as the
source observes \cite[487]{Keisler1961ElementaryClasses}.  The source also
notes that for structures with no relations, and for finite structures, the
equivalence was known without that instance
\cite[486, footnote~4]{Keisler1961ElementaryClasses}.

\begin{hbtheorem}{thm:p03-w03-p03-018-keisler-gch}{proved=annex:p03-w03-p03-018-keisler-gch}{Keisler's existential form}
Assume the generalized continuum hypothesis.  Then
\(\mathfrak A^*=\mathfrak B^*\) if and only if
\(\mathrm{Pow}(\mathfrak A)\cap\mathrm{Pow}(\mathfrak B)\ne\varnothing\).
\end{hbtheorem}

This is Theorem~2.4 of the source \cite[487]{Keisler1961ElementaryClasses}, stated under a
standing assumption of the generalized continuum hypothesis for the rest of
the source's Section~2.  Its forward direction uses that hypothesis at one
cardinal \(\alpha\) only, chosen large enough for the two structures and their
type.  Keisler adds that, without the hypothesis, it was not known whether the
relation \(\mathrm{Pow}(\mathfrak A)\cap\mathrm{Pow}(\mathfrak B)\ne\varnothing\)
is transitive \cite[487]{Keisler1961ElementaryClasses}, and he closes
Section~2 by asking whether his Theorems~2.4--2.14 can be proved without it
\cite[490]{Keisler1961ElementaryClasses}.

The appendix of the source announces, without proofs
\cite[490]{Keisler1961ElementaryClasses}, a stronger form in which the
ultrafilters are fixed before the structures.  It uses the following class of
ultrafilters \cite[491]{Keisler1961ElementaryClasses}.  A filter \(D\) on
\(I\) is \emph{countably complete} when \(\bigcap X\in D\) for every
\(X\subseteq D\) with \(\varkappa(X)\le\omega\).  For a cardinal \(\alpha\),
\(D\in Q(\alpha)\) means: \(D\) is an ultrafilter on some set \(I\) (so
\(I=\bigcup D\)); \(D\) is not countably complete; and for every cardinal
\(\beta<\alpha\) and every function \(F\) from \(\overline S_\omega(\beta)\)
into \(D\) that is monotonic (\(t\subseteq u\) implies
\(F(t)\subseteq F(u)\)), there is a function \(G\) from
\(\overline S_\omega(\beta)\) into \(D\) that is multiplicative
(\(G(t\cap u)=G(t)\cap G(u)\)) with \(G(t)\subseteq F(t)\) for every
\(t\in\overline S_\omega(\beta)\).  From that point on the source supposes
\(\alpha,\beta\ge\omega\).

\begin{hbtheorem}{thm:p03-w03-p03-018-keisler-appendix}{cited=Keisler1961ElementaryClasses}{Keisler's announced ultrafilter-first forms}
\begin{enumerate}
  \item \(\alpha^+=\varkappa(2^\alpha)\) if and only if there is an
        ultrafilter \(D\) on \(\alpha\) with \(D\in Q(\varkappa(2^\alpha))\).
  \item Suppose \(\beta^+=\varkappa(2^\beta)\), \(\beta\ge\varkappa(\varrho)\),
        \(D,E\in Q(\beta^+)\), \(\bigcup D=\bigcup E=\beta\), and
        \(\mathfrak A\), \(\mathfrak B\) are of cardinality at most
        \(\beta^+\).  Then \(\mathfrak A^*=\mathfrak B^*\) if and only if
        \(\mathfrak A^\beta/D\cong\mathfrak B^\beta/E\).
  \item \(\mathfrak A^*=\mathfrak B^*\) if and only if
        \(\mathfrak A^{\bigcup D}/D\cong\mathfrak B^{\bigcup D}/D\) for some
        ultrafilter \(D\).
\end{enumerate}
\end{hbtheorem}

These are item~(iii) of Lemma~A.3 \cite[491]{Keisler1961ElementaryClasses}
and Corollaries~A.6 and~A.8 \cite[492]{Keisler1961ElementaryClasses}.  The
source gives no proof of them, and none is given here; they stand on the
source's announcement alone.  Corollary~A.8 is introduced as an improved form
of Theorem~2.4 and is printed with no hypothesis in its statement; the
appendix states the generalized continuum hypothesis as a standing assumption
only later, before its results on reduced products
\cite[493]{Keisler1961ElementaryClasses}.

\subsection{Keisler's and Shelah's hypotheses compared}

\begin{hbtheorem}{thm:p03-w03-p03-018-keisler-compare}{proved=annex:p03-w03-p03-018-keisler-compare}{Comparison of the two theorems' hypotheses}
\begin{enumerate}
  \item Let \(\varkappa(I)=\alpha\ge\omega\) and let \(D\) be an ultrafilter
        on \(I\).  Then \(\overline S_\alpha(I)\subseteq D\) if and only if
        every member of \(D\) has cardinality \(\alpha\).  Such a \(D\) is
        not principal.
  \item Let \(\alpha\ge\omega\) with \(\varkappa(2^\alpha)=\alpha^+\), and let
        \(\mu\) be the bound of \cref{def:p03-w03-p03-018-notation} for
        \(\lambda=\alpha\).  Then \(\mu\le\alpha<\varkappa(2^\alpha)\).  Every
        structure of power less than \(\mu\) has cardinality at most
        \(\varkappa(2^\alpha)\), while the structure with universe \(\alpha\)
        and no relations has cardinality at most \(\varkappa(2^\alpha)\) and
        power not less than \(\mu\).
  \item With no cardinal-arithmetic hypothesis beyond the ambient set theory
        with Choice, for structures of one type,
        \(\mathrm{Pow}(\mathfrak A)\cap\mathrm{Pow}(\mathfrak B)\ne\varnothing\)
        holds if and only if \(\mathfrak A^*=\mathfrak B^*\).  In particular
        this relation is an equivalence relation.
\end{enumerate}
\end{hbtheorem}

Item~3 settles the transitivity question that Keisler left open; it rests on
Shelah's construction as proved in \cref{thm:p03-w03-p03-017-keisler-shelah}.
The two theorems compare along four lines.

\begin{itemize}
  \item \emph{Cardinal arithmetic.}
        \Cref{thm:p03-w03-p03-018-keisler-products,thm:p03-w03-p03-018-keisler-powers}
        assume \(\varkappa(2^\alpha)=\alpha^+\) at the cardinal \(\alpha\)
        that is also the cardinality of the index set, and
        \cref{thm:p03-w03-p03-018-keisler-gch} assumes the generalized
        continuum hypothesis.  \Cref{thm:p03-w03-p03-018-shelah} assumes
        nothing beyond the ambient set theory with Choice
        \cite[224]{Shelah1971Ultrapowers}.
  \item \emph{Index set and size of the structures.}  Both index by the
        governing cardinal: \(\alpha\) for Keisler, \(\lambda\) for Shelah.
        At a common cardinal \(\alpha=\lambda\) where Keisler's hypothesis
        holds, Keisler covers structures of cardinality up to \(\alpha^+\),
        while Shelah covers only power below \(\mu\le\alpha\) (item~2).
        Shelah pays for removing the arithmetic hypothesis with a smaller
        bound, so a given pair needs a larger \(\lambda\)
        (\cref{thm:p03-w03-p03-018-index-bound}).
  \item \emph{Language.}  Keisler requires a relational language with at most
        \(\alpha\) predicate symbols.  Shelah's statement puts no bound on the
        language; his proof reduces to a language of cardinality at most
        \(\lambda\) \cite[230]{Shelah1971Ultrapowers}.
  \item \emph{Ultrafilters.}  Keisler's proved theorems give two ultrafilters
        \(D\) and \(E\), possibly different, constructed after the structures
        are given; each contains \(\overline S_\alpha(\alpha)\) and so is
        nonprincipal (item~1).  Shelah's theorem gives one ultrafilter \(D\)
        over \(\lambda\), fixed before any structure is mentioned, that serves
        every pair below the bound.  A single ultrafilter fixed in advance
        appears in Keisler's article only in the announced, unproved
        \cref{thm:p03-w03-p03-018-keisler-appendix}.
\end{itemize}

Shelah summarizes Keisler's theorem as giving, when \(\lambda^+=2^\lambda\),
one ultrafilter \(D\) over \(\lambda\) with
\(M^\lambda/D\cong N^\lambda/D\) for all \(M\equiv N\) of power at most
\(\lambda^+\) in a language of cardinality at most \(\lambda\)
\cite[224]{Shelah1971Ultrapowers}; his reference is to the article read here
\cite[232]{Shelah1971Ultrapowers}.  Read against the article, that summary is
item~2 of \cref{thm:p03-w03-p03-018-keisler-appendix} with \(\beta=\lambda\)
and \(D=E\), the existence of such a \(D\) on \(\lambda\) coming from item~1;
both are announced there without proof.  The theorems the article proves,
\cref{thm:p03-w03-p03-018-keisler-products,thm:p03-w03-p03-018-keisler-powers},
have the same cardinal hypotheses but two ultrafilters chosen for the given
pair.

\begin{hbexample}{def:p03-w03-p03-018-keisler-example}{The rationals and the reals under the continuum hypothesis}
Assume \(\varkappa(2^\omega)=\omega^+\), and let
\(\mathfrak A=(\mathbb Q,<)\) and \(\mathfrak B=(\mathbb R,<)\).  Both are
dense strict linear orders without endpoints (for \(x<y\) the points
\(x-1\), \((x+y)/2\) and \(y+1\) witness this), so
\(\mathfrak A^*=\mathfrak B^*\) by \cref{thm:p03-rn-dlo-pattern-truth}.  The
type has one predicate, so \(\varkappa(\varrho)=1\le\omega\), and both
structures have cardinality at most \(\varkappa(2^\omega)\).
\Cref{thm:p03-w03-p03-018-keisler-powers} with \(\alpha=\omega\) gives
ultrafilters \(D\), \(E\) on \(\omega\), each containing every cofinite
subset of \(\omega\), with
\(\mathfrak A^\omega/D\cong\mathfrak B^\omega/E\).

Shelah's theorem cannot do this at \(\lambda=\omega\): there
\(\mu=\aleph_0\) (\cref{thm:p03-w03-p03-018-mu}), and only finite models are
covered.  Nor can it at \(\lambda=2^{\aleph_0}\), since \(\mathbb R\) has
power \(2^{\aleph_0}\), which is not below \(\mu\) by the computation in
\cref{def:p03-w03-p03-018-worked}.  At \(\lambda=2^{2^{\aleph_0}}\) it gives
one ultrafilter \(D\) with
\(\mathfrak A^\lambda/D\cong\mathfrak B^\lambda/D\), with or without the
continuum hypothesis. These applications give two sufficient index bounds:
Keisler's countable one under \(\varkappa(2^\omega)=\omega^+\), and
Shelah's displayed larger one without that hypothesis. The comparison does
not establish that the larger bound is necessary.
\end{hbexample}

\begin{hbwarning}{def:p03-w03-p03-018-keisler-limits}{Limits of Keisler's version}
\begin{itemize}
  \item Keisler's proof needs \(\varkappa(2^\alpha)=\alpha^+\) at the chosen
        \(\alpha\): it is what bounds the products \(\prod_iA_i\) and
        \(\prod_iB_i\) by \(\alpha^+\), so that they can be exhausted in
        \(\alpha^+\) steps, each step handling fewer than \(\alpha^+\)
        named elements.
  \item The two ultrafilters of \cref{thm:p03-w03-p03-018-keisler-powers}
        depend on the pair, and nothing in the proof makes \(D=E\).  That one
        ultrafilter fixed in advance suffices is, in the source, only the
        announced \cref{thm:p03-w03-p03-018-keisler-appendix}.
  \item The source reports that a stronger result announced by Kochen, from
        which Kochen derived the existential form, is false, so that the
        proof Kochen had in mind does not seem correct
        \cite[487, footnote~5]{Keisler1961ElementaryClasses}.
  \item As with Shelah's theorem, isomorphic ultrapowers do not make the
        structures isomorphic: in the example above \(\mathbb Q\) and
        \(\mathbb R\) have different cardinalities.
\end{itemize}
\end{hbwarning}

\subsection{Recovering the existential form}

\begin{hbcorollary}{thm:p03-w03-p03-018-index-bound}{proved=annex:p03-w03-p03-018-index-bound}{An explicit index set for elementarily equivalent models}
Let \(M\equiv N\) and put
\(\kappa=\lVert M\rVert+\lVert N\rVert+\aleph_0\) and \(\lambda=2^\kappa\).
Then there is an ultrafilter \(D\) over \(\lambda\) with
\(M^\lambda/D\cong N^\lambda/D\).  Consequently the forward implication of
\cref{thm:p03-keisler-shelah} holds with index set \(2^\kappa\) for both
ultrapowers and with one ultrafilter for both.
\end{hbcorollary}

Part~1 of \cref{thm:p03-w03-p03-018-shelah}, used in this corollary,
is proved in \cref{thm:p03-w03-p03-017-shelah-proof}. Thus the corollary
depends on that construction and the cardinal computation above, without
using the still-cited part~3 or the goodness assertion. The reverse
implication of \cref{thm:p03-keisler-shelah} needs no cardinal hypothesis at
all; it is item~2 of the next result.

\subsection{Equivalence, isomorphic ultrapowers and isomorphism}

The three relations below are distinct for infinite structures.  Isomorphism
of \(M\) and \(N\) is the strongest; isomorphism of some ultrapowers is
intermediate in formulation; and elementary equivalence is the weakest in
formulation.  \Cref{thm:p03-w03-p03-018-shelah} shows that the last two
coincide; the first is genuinely stronger.

\begin{hbtheorem}{thm:p03-w03-p03-018-comparison}{proved=annex:p03-w03-p03-018-comparison}{Comparison of the three relations}
Let \(M\) and \(N\) be structures of one signature, let \(I\) and \(J\) be
nonempty sets, and let \(D\) and \(E\) be ultrafilters on \(I\) and \(J\).
\begin{enumerate}
  \item If \(M\cong N\), then \(M^I/D\cong N^I/D\).
  \item If \(M^I/D\cong N^J/E\), then \(M\equiv N\).
  \item If \(D\) is principal, generated by \(\{k\}\) for some \(k\in I\),
        then \(M^I/D\cong M\).
  \item If \(M\) is finite and \(M\equiv N\), then \(M\cong N\).  Hence, for
        finite \(M\), the conditions \(M\equiv N\), \(M\cong N\) and
        \(M^I/D\cong N^I/D\) are equivalent.
\end{enumerate}
\end{hbtheorem}

For infinite structures item~4 fails, and an ultrapower isomorphism does not
reflect back to the factors.

\begin{hbcorollary}{thm:p03-w03-p03-018-separation}{proved=annex:p03-w03-p03-018-separation}{Isomorphic ultrapowers of non-isomorphic models}
Order \(2^\omega\), the set of infinite \(0\)--\(1\) sequences,
lexicographically: \(s<_{\mathrm{lex}}t\) when \(s\ne t\) and
\(s(n)<t(n)\) at the least \(n\) with \(s(n)\ne t(n)\).  Order
\(2^\omega\times\mathbb Q\) by \((s,q)<(t,r)\) when \(s<_{\mathrm{lex}}t\),
or \(s=t\) and \(q<r\).  Let \(M=(\mathbb Q,<)\) and
\(N=(2^\omega\times\mathbb Q,<)\), and let \(\lambda=2^{2^{\aleph_0}}\).
Then:
\begin{enumerate}
  \item \(M\equiv N\);
  \item \(M\not\cong N\);
  \item for the ultrafilter \(D\) over \(\lambda\) given by
        \cref{thm:p03-w03-p03-018-shelah}, \(M^\lambda/D\cong N^\lambda/D\);
  \item that \(D\) is not principal; and
  \item \(M\not\cong M^\lambda/D\), although \(M\equiv M^\lambda/D\); the
        diagonal image of \(M\) is a proper elementary substructure of
        \(M^\lambda/D\).
\end{enumerate}
\end{hbcorollary}

\begin{hbexample}{def:p03-w03-p03-018-worked}{Choosing \(\lambda\) for a given pair}
Take \(M=(\mathbb Q,<)\) and \(N\) as in
\cref{thm:p03-w03-p03-018-separation}, of powers \(\aleph_0\) and
\(2^{\aleph_0}\).  The choice \(\lambda=\aleph_0\) is useless: then
\(\mu=\aleph_0\) by \cref{thm:p03-w03-p03-018-mu}, and part~1 of the theorem
covers only finite models, for which item~4 of
\cref{thm:p03-w03-p03-018-comparison} already gives isomorphism.  The choice
\(\lambda=2^{\aleph_0}\) gives \(\mu>\aleph_0\), so it covers \(M\) together
with any countable dense linear order without endpoints, but not \(N\)
unless \(2^{\aleph_0}<\mu\), which would require
\((2^{\aleph_0})^{2^{\aleph_0}}=2^{\aleph_0}\), contradicting
\(2^{2^{\aleph_0}}>2^{\aleph_0}\).  The choice
\(\lambda=2^{2^{\aleph_0}}\) gives \(\mu>2^{\aleph_0}\), so both models fall
under part~1, and one ultrafilter \(D\) over \(2^{2^{\aleph_0}}\) yields
\(M^\lambda/D\cong N^\lambda/D\).  This is the computation behind
\cref{thm:p03-w03-p03-018-index-bound} with \(\kappa=2^{\aleph_0}\).
\end{hbexample}

\begin{hbwarning}{def:p03-w03-p03-018-limits}{What the theorem does not give}
\begin{itemize}
  \item It does not make \(M\) and \(N\) isomorphic, nor make \(M\)
        isomorphic to its own ultrapower:
        \cref{thm:p03-w03-p03-018-separation} has elementarily equivalent,
        non-isomorphic models with isomorphic ultrapowers, and a model not
        isomorphic to its ultrapower.
  \item It does not describe \(D\).  The ultrafilter is obtained by a
        transfinite construction in the ambient set theory with Choice, and
        a principal ultrafilter never works for a non-isomorphic pair.
  \item It is not uniform in \(\lambda\): models of power at least \(\mu\)
        are not covered for that \(\lambda\), and \(\mu\le\lambda\) always.
        A larger pair needs a larger \(\lambda\), as in
        \cref{thm:p03-w03-p03-018-index-bound}.
  \item Its first-order hypothesis cannot be dropped: by item~2 of
        \cref{thm:p03-w03-p03-018-comparison}, models that are not
        elementarily equivalent have no isomorphic ultrapowers for any index
        sets or ultrafilters.
\end{itemize}
\end{hbwarning}

\paragraph{Relationships.}
Item~2 of \cref{thm:p03-w03-p03-018-comparison} is \L o\'s's theorem
(\cref{thm:los}) through the diagonal embedding
(\cref{cor:diagonal-ultrapower}) and isomorphism invariance
(\cref{lem:iso-invariance}).  The completeness of the theory of dense linear
orders without endpoints used in \cref{thm:p03-w03-p03-018-separation} is
\cref{thm:p03-rn-dlo-pattern-truth}.  The use of isomorphic ultrapowers in
the joint-consistency argument of the Part~III annex relies on
\cref{thm:p03-keisler-shelah}, whose index sets are made explicit by
\cref{thm:p03-w03-p03-018-index-bound}.  The Choice assumptions are recorded
in \cref{def:p03-ultrafilter-choice-boundary}.

\paragraph{Sources.}
Shelah's theorem and its hypotheses are from S.~Shelah, \emph{Every two
elementarily equivalent models have isomorphic ultrapowers}, Israel Journal
of Mathematics 10 (1971), pages 224--233 \cite{Shelah1971Ultrapowers}: the
statement and the summary of Keisler's result on page 224, the notation on
page 225, the reduction of the language and the single list of model pairs on
page 230, and the reference to Keisler's article on page 232.  Keisler's
theorems are from H.~J.~Keisler, \emph{Ultraproducts and elementary
classes}, Indagationes Mathematicae 23 (1961), pages 477--495
\cite{Keisler1961ElementaryClasses}: the standing use of Choice on page 479,
the set-theoretic notation and the generalized continuum hypothesis on
page 480, the structures and language on pages 480--481, ultraproducts and
\(\mathrm{Pow}\) on page 483, Lemma~2.1 and Theorem~2.2 on page 484 with
the proof on pages 484--486, Theorem~2.3 on page 486, Theorem~2.4 and the
transitivity remark on page 487, the open question on page 490, and the
appendix results on pages 491--492.  Both presentations follow the published
journal pages.
